\subsection{Graph Coloring}
\begin{frame}{\(m\)-Coloring의 State}
\begin{description}
\item[State] 정해진 order의 first \(k\) vertices가 colored.
\item[Decision] vertex \(k\)에 color \(1..m\).
\item[Promising] 이미 colored인 neighbor와 다른 color.
\item[Complete] 모든 vertex가 colored.
\end{description}
\end{frame}
\begin{frame}{작은 Graph Coloring Trace}
\centering\begin{tikzpicture}
\node[gvertex](a)at(-2,1){A};\node[gvertex](b)at(0,2){B};\node[gvertex](c)at(2,1){C};\node[gvertex](d)at(0,-.5){D};
\foreach \u/\v in{a/b,b/c,c/d,d/a,a/c}\draw[gedge](\u)--(\v);
\only<1|handout:0>{\node[ssnote]at(4,-.5){order: A, B, C, D --- 아직 색이 없다};}
\visible<2->{\node[gvertex,fill=red!20](aa)at(-2,1){A:1};}
\only<2|handout:0>{\node[ssnote]at(4,-.5){A: 색 1 --- 이웃이 아직 없다};}
\visible<3->{\node[gvertex,fill=AlgoBlue!18](bb)at(0,2){B:2};\node[font=\tiny,text=red!60!black,right=1mm of bb]{\(\times\)1};}
\only<3|handout:0>{\node[ssnote]at(4,-.5){B: 색 1은 A와 충돌 \(\to\) 색 2};}
\visible<4->{\node[gvertex,fill=green!15](cc)at(2,1){C:3};\node[font=\tiny,text=red!60!black,right=1mm of cc]{\(\times\)1 \(\times\)2};}
\only<4|handout:0>{\node[ssnote]at(4,-.5){C: 색 1은 A와, 색 2는 B와 충돌 \(\to\) 색 3};}
\visible<5->{\node[gvertex,fill=AlgoBlue!18](dd)at(0,-.5){D:2};\node[font=\tiny,text=red!60!black,right=1mm of dd]{\(\times\)1};}
\only<5|handout:1>{\node[ssinc]at(4,-.5){valid 3-coloring};\node[font=\scriptsize]at(4,-1.3){D의 이웃은 A,C뿐 --- B=2지만 이웃 아님};}
\end{tikzpicture}
\end{frame}
\begin{frame}{Coloring Pseudocode}
\algofont\begin{algorithmic}[1]
\Procedure{Color}{$k$}
\If{\(k=|V|\)}\State output color; \Return true\EndIf
\State \(v\gets order[k]\)
\For{\(c\gets1\) to \(m\)}
\If{every colored \(u\in Adj[v]\) has \(color[u]\ne c\)}
\State \(color[v]\gets c\)
\If{\Call{Color}{k+1}}\State\Return true\EndIf
\State \(color[v]\gets0\)
\EndIf\EndFor
\State\Return false
\EndProcedure
\end{algorithmic}
\end{frame}
\begin{frame}{Vertex Order가 성능을 바꾼다}
\begin{itemize}
\item 높은 degree 또는 제약이 강한 vertex를 먼저 고르면 conflict를 일찍 발견할 수 있다.
\item order가 바뀌어도 모든 color를 시도하면 completeness는 유지된다.
\item 동적 “most constrained variable”은 CSP heuristic의 예다.
\end{itemize}
\end{frame}
\begin{frame}{Coloring Correctness}
\begin{enumerate}
\item 각 level은 한 vertex의 모든 color를 열거한다.
\item 이미 같은 색 neighbor가 있으면 그 conflict는 descendant에서도 사라지지 않는다.
\item 따라서 prune은 sound하다.
\end{enumerate}
\end{frame}
\begin{frame}{Graph Coloring Takeaway}
\sstakeaway{Graph Coloring도 같은 Backtracking skeleton을 사용한다. 달라지는 것은 state, Candidates, Promising뿐이다.}
\end{frame}
